%----------------------------------------------------------------
%
%  File    :  survey-relatedwork.tex
%
%  Author  :   Lukas Auer, David Heidinger, Nina Tschikof and Christina Vogel, TU Graz, Austria
% 
%  Created :  06 May 2026
% 
%  Changed :  06 May 2026
% 
%----------------------------------------------------------------


\chapter{Related Work}
\label{chap:RelatedWork}

This chapter presents prior work relevant to the identification of the nine dishonest chart techniques described in chapter \ref{chap:Taxonomy}. The reviewed sources include academic papers, books, and online resources.

\section{Lan and Liu}
"I Came Across A Junk" is a paper by Xingyu Lan and Yu Liu from 2025 \cite{Liu2025}. To understand design flaws of data visualisation, the authors analysed 2227 flawed data visualisations from which they derived a taxonomy of 76 specific design flaws. These design flaws were categorised into three high-level categories and ten subcategories. They then identified seven causes of the flaws and proposed a research agenda for combating these flaws as well as present future research opportunities in this field.

The misinformation category includes flaws that distort or misrepresent information, such as truncated axes, misleading 3D effects, or incompatible chart types. Uninformativeness describes visualisations that either overwhelm users with excessive complexity or omit important contextual information. Unsociability focuses on visualisations that create discomfort or confusion through abnormal or aggressive design choices.
\\\\
The authors conducted focus group studies to find out why these design flaws are happening. The authors identified seven major causes of these design flaws: Achieving communication intents; Lack of literacy and expertise; Pursuit of aesthetics; Data that is difficult to handle; Limitations of tools; Compromise with clients; Ineffective team collaboration.
\\\\
The authors also published a website called FlawViz where the taxonomy, the data corpus used for the study, as well as the paper can be viewed.
\\\\
The taxonomy proposed by Lan and Liu is particularly relevant for this work because it provides a broad and structured overview of misleading and flawed visualisation techniques. Several dishonest chart techniques discussed in chapter \ref{chap:Taxonomy} are directly related to flaws classified under misinformation, especially inaccuracy and unfairness.


\section{Cairo}
The book "How Charts Lie" by Alberto Cairo offers an in-depth analysis of how charts can be used for either revealing, or concealing information \cite{cairo2019}. The book explains the fundamental principles of chart design and why they are an important tool in making trends in data visible. Then multiple chapters of different ways in which charts can be used to lie to the viewer are presented. Cairo argues that charts are often perceived as objective and trustworthy by the general public, which makes misleading visualisations particularly influential in shaping opinions and public understanding. It is therefore important to be truthful and intentional when designing a chart.
\\\\
The book uses many real world examples of misleading design choices and how they can distort the perception of data. These include issues such as truncated axes, inappropriate chart types, missing context, and misleading proportions. 
\\\\
For this survey, the book was particularly relevant because it provides both practical examples and conceptual explanations of deceptive visualisation techniques. Especially the early chapters helped identify common patterns of misleading chart design and supported the development of the taxonomy presented in chapter \ref{chap:Taxonomy}.n

 
\section{Nathan Yau}
The online guide "Defense Against Dishonest Charts" by Nathan Yau, published on the FlowingData website presents interactive examples of misleading chart design and explains how visualisations can distort the interpretation of data \cite{Yau2025}. By allowing users to interactively compare misleading and corrected versions of charts, the website highlights how small design decisions can significantly influence the perception of information.
\\\\
This resource was particularly useful for this survey because it illustrates the visual impact of misleading design choices in a concise and practical way. Furthermore, the effectiveness of the interactive examples inspired the development of interactive charts for our proposed taxonomy of misleading chart techniques.


\section{UKPLab}
The GitHub repository "Awesome Misleading Visualizations", maintained by the UKP Lab at TU Darmstadt, provides a curated collection of papers, datasets, and online resources related to misleading data visualisations and chart understanding \cite{UKPLab2025}. 
\\\\
This resource was useful for this survey because it served as a central entry point for finding literature and datasets in the field of misleading visualisations.

\section{Wijnker et al.}
The paper "Debunking Strategies for Misleading Bar Charts" by Wijnker et al., investigates how misleading bar charts can be corrected and how viewers can be protected from misinformation caused by deceptive visualisations \cite{Wijnker2022}. The authors focus specifically on bar charts with truncated y-axes.
The paper presents an experimental study involving two online surveys with more than 400 participants. The findings show that the proposed correction methods reduced the misleading effect of manipulated charts, while presenting a corrected alternative chart was found to be the most effective strategy.
\\\\
We did not use this paper for our presentation but it was helpful in gathering insights and information about the topic of misleading chart techniques.


\section{Lea Gaslowitz}
The animated video "How to Spot a Misleading Graph" by Lea Gaslowitz presented on TED-Ed explains how axis manipulation and missing context distorts data \cite{Gaslowitz2017}. 
\\\\
We did not use this paper for our presentation but it was helpful in gathering insights and information about the topic of misleading chart techniques.
